% Building a Sun RPC client and server with rpcgen from sum.x.
% White: written by the developer; shaded: generated by rpcgen.
% Usage: % Building a Sun RPC client and server with rpcgen from sum.x.
% White: written by the developer; shaded: generated by rpcgen.
% Usage: % Building a Sun RPC client and server with rpcgen from sum.x.
% White: written by the developer; shaded: generated by rpcgen.
% Usage: % Building a Sun RPC client and server with rpcgen from sum.x.
% White: written by the developer; shaded: generated by rpcgen.
% Usage: \input{figures/l13-rpcgen}   (width ~120 mm)
\begin{tikzpicture}[x=1mm, y=1mm, font=\footnotesize\sffamily,
  >={Stealth[length=1.5mm]},
  file/.style={draw, fill=white, align=center, minimum width=18.5mm,
               minimum height=9mm, inner sep=0.5pt},
  gen/.style={file, fill=ln-accent!15},
  tool/.style={draw, rounded corners=3pt, fill=ln-accent!40, align=center,
               minimum width=22mm, minimum height=6.5mm},
  exe/.style={draw, thick, rounded corners=2pt, fill=ln-box, align=center,
              minimum width=30mm, minimum height=7mm},
  fn/.style={font=\scriptsize\ttfamily},
  what/.style={font=\tiny\sffamily, text=ln-muted},
  lab/.style={font=\scriptsize\sffamily, text=ln-muted}]
  % English descriptions at \scriptsize: \tiny is unreadable on screen, and
  % the widest ("types, constants", 16.4 mm) still fits the 18.5 mm box.
  \newcommand{\lnFile}[2]{{\scriptsize\ttfamily #1}\\[-0.4ex]{\iflangja{\tiny}{\scriptsize}\sffamily\color{ln-muted}#2}}
  % interface and compiler
  \node[file] (x) at (60,52) {\lnFile{sum.x}{\iflangja{XDR インタフェース}{XDR interface}}};
  \node[tool] (rpcgen) at (60,40) {\texttt{rpcgen}};
  \draw[->] (x) -- (rpcgen);
  % source files
  \node[file] (cli) at (10,24) {\lnFile{client.c}{\iflangja{呼び出し元}{caller}}};
  \node[gen]  (cln) at (30,24) {\lnFile{sum\_clnt.c}{\iflangja{クライアントスタブ}{client stub}}};
  \node[gen]  (h)   at (50,24) {\lnFile{sum.h}{\iflangja{型と定数}{types, constants}}};
  \node[gen]  (xdr) at (70,24) {\lnFile{sum\_xdr.c}{\iflangja{XDR ルーチン}{XDR routines}}};
  \node[gen]  (svc) at (90,24) {\lnFile{sum\_svc.c}{\iflangja{サーバスタブ}{server stub}}};
  \node[file] (srv) at (110,24) {\lnFile{server.c}{\iflangja{手続きの実装}{procedure}}};
  \foreach \f in {cln, h, xdr, svc} \draw[->] (rpcgen.south) -- (\f.north);
  % legend
  \node[file, minimum width=5mm, minimum height=3.5mm] (l1) at (3,53) {};
  \node[lab, anchor=west] at (6,53) {\iflangja{開発者が記述}{written by the developer}};
  \node[gen, minimum width=5mm, minimum height=3.5mm] (l2) at (3,48) {};
  \node[lab, anchor=west] at (6,48) {\iflangja{\texttt{rpcgen} が生成}{generated by \texttt{rpcgen}}};
  % executables
  \node[exe] (ce) at (30,5) {\iflangja{クライアントプログラム}{client program}};
  \node[exe] (se) at (90,5) {\iflangja{サーバプログラム}{server program}};
  \foreach \f in {cli, cln, xdr} \draw[->] (\f.south) -- (ce.north);
  \foreach \f in {xdr, svc, srv} \draw[->] (\f.south) -- (se.north);
  \node[lab] at (60,12.5) {\iflangja{コンパイル・リンク}{compile and link}};
\end{tikzpicture}
   (width ~120 mm)
\begin{tikzpicture}[x=1mm, y=1mm, font=\footnotesize\sffamily,
  >={Stealth[length=1.5mm]},
  file/.style={draw, fill=white, align=center, minimum width=18.5mm,
               minimum height=9mm, inner sep=0.5pt},
  gen/.style={file, fill=ln-accent!15},
  tool/.style={draw, rounded corners=3pt, fill=ln-accent!40, align=center,
               minimum width=22mm, minimum height=6.5mm},
  exe/.style={draw, thick, rounded corners=2pt, fill=ln-box, align=center,
              minimum width=30mm, minimum height=7mm},
  fn/.style={font=\scriptsize\ttfamily},
  what/.style={font=\tiny\sffamily, text=ln-muted},
  lab/.style={font=\scriptsize\sffamily, text=ln-muted}]
  % English descriptions at \scriptsize: \tiny is unreadable on screen, and
  % the widest ("types, constants", 16.4 mm) still fits the 18.5 mm box.
  \newcommand{\lnFile}[2]{{\scriptsize\ttfamily #1}\\[-0.4ex]{\iflangja{\tiny}{\scriptsize}\sffamily\color{ln-muted}#2}}
  % interface and compiler
  \node[file] (x) at (60,52) {\lnFile{sum.x}{\iflangja{XDR インタフェース}{XDR interface}}};
  \node[tool] (rpcgen) at (60,40) {\texttt{rpcgen}};
  \draw[->] (x) -- (rpcgen);
  % source files
  \node[file] (cli) at (10,24) {\lnFile{client.c}{\iflangja{呼び出し元}{caller}}};
  \node[gen]  (cln) at (30,24) {\lnFile{sum\_clnt.c}{\iflangja{クライアントスタブ}{client stub}}};
  \node[gen]  (h)   at (50,24) {\lnFile{sum.h}{\iflangja{型と定数}{types, constants}}};
  \node[gen]  (xdr) at (70,24) {\lnFile{sum\_xdr.c}{\iflangja{XDR ルーチン}{XDR routines}}};
  \node[gen]  (svc) at (90,24) {\lnFile{sum\_svc.c}{\iflangja{サーバスタブ}{server stub}}};
  \node[file] (srv) at (110,24) {\lnFile{server.c}{\iflangja{手続きの実装}{procedure}}};
  \foreach \f in {cln, h, xdr, svc} \draw[->] (rpcgen.south) -- (\f.north);
  % legend
  \node[file, minimum width=5mm, minimum height=3.5mm] (l1) at (3,53) {};
  \node[lab, anchor=west] at (6,53) {\iflangja{開発者が記述}{written by the developer}};
  \node[gen, minimum width=5mm, minimum height=3.5mm] (l2) at (3,48) {};
  \node[lab, anchor=west] at (6,48) {\iflangja{\texttt{rpcgen} が生成}{generated by \texttt{rpcgen}}};
  % executables
  \node[exe] (ce) at (30,5) {\iflangja{クライアントプログラム}{client program}};
  \node[exe] (se) at (90,5) {\iflangja{サーバプログラム}{server program}};
  \foreach \f in {cli, cln, xdr} \draw[->] (\f.south) -- (ce.north);
  \foreach \f in {xdr, svc, srv} \draw[->] (\f.south) -- (se.north);
  \node[lab] at (60,12.5) {\iflangja{コンパイル・リンク}{compile and link}};
\end{tikzpicture}
   (width ~120 mm)
\begin{tikzpicture}[x=1mm, y=1mm, font=\footnotesize\sffamily,
  >={Stealth[length=1.5mm]},
  file/.style={draw, fill=white, align=center, minimum width=18.5mm,
               minimum height=9mm, inner sep=0.5pt},
  gen/.style={file, fill=ln-accent!15},
  tool/.style={draw, rounded corners=3pt, fill=ln-accent!40, align=center,
               minimum width=22mm, minimum height=6.5mm},
  exe/.style={draw, thick, rounded corners=2pt, fill=ln-box, align=center,
              minimum width=30mm, minimum height=7mm},
  fn/.style={font=\scriptsize\ttfamily},
  what/.style={font=\tiny\sffamily, text=ln-muted},
  lab/.style={font=\scriptsize\sffamily, text=ln-muted}]
  % English descriptions at \scriptsize: \tiny is unreadable on screen, and
  % the widest ("types, constants", 16.4 mm) still fits the 18.5 mm box.
  \newcommand{\lnFile}[2]{{\scriptsize\ttfamily #1}\\[-0.4ex]{\iflangja{\tiny}{\scriptsize}\sffamily\color{ln-muted}#2}}
  % interface and compiler
  \node[file] (x) at (60,52) {\lnFile{sum.x}{\iflangja{XDR インタフェース}{XDR interface}}};
  \node[tool] (rpcgen) at (60,40) {\texttt{rpcgen}};
  \draw[->] (x) -- (rpcgen);
  % source files
  \node[file] (cli) at (10,24) {\lnFile{client.c}{\iflangja{呼び出し元}{caller}}};
  \node[gen]  (cln) at (30,24) {\lnFile{sum\_clnt.c}{\iflangja{クライアントスタブ}{client stub}}};
  \node[gen]  (h)   at (50,24) {\lnFile{sum.h}{\iflangja{型と定数}{types, constants}}};
  \node[gen]  (xdr) at (70,24) {\lnFile{sum\_xdr.c}{\iflangja{XDR ルーチン}{XDR routines}}};
  \node[gen]  (svc) at (90,24) {\lnFile{sum\_svc.c}{\iflangja{サーバスタブ}{server stub}}};
  \node[file] (srv) at (110,24) {\lnFile{server.c}{\iflangja{手続きの実装}{procedure}}};
  \foreach \f in {cln, h, xdr, svc} \draw[->] (rpcgen.south) -- (\f.north);
  % legend
  \node[file, minimum width=5mm, minimum height=3.5mm] (l1) at (3,53) {};
  \node[lab, anchor=west] at (6,53) {\iflangja{開発者が記述}{written by the developer}};
  \node[gen, minimum width=5mm, minimum height=3.5mm] (l2) at (3,48) {};
  \node[lab, anchor=west] at (6,48) {\iflangja{\texttt{rpcgen} が生成}{generated by \texttt{rpcgen}}};
  % executables
  \node[exe] (ce) at (30,5) {\iflangja{クライアントプログラム}{client program}};
  \node[exe] (se) at (90,5) {\iflangja{サーバプログラム}{server program}};
  \foreach \f in {cli, cln, xdr} \draw[->] (\f.south) -- (ce.north);
  \foreach \f in {xdr, svc, srv} \draw[->] (\f.south) -- (se.north);
  \node[lab] at (60,12.5) {\iflangja{コンパイル・リンク}{compile and link}};
\end{tikzpicture}
   (width ~120 mm)
\begin{tikzpicture}[x=1mm, y=1mm, font=\footnotesize\sffamily,
  >={Stealth[length=1.5mm]},
  file/.style={draw, fill=white, align=center, minimum width=18.5mm,
               minimum height=9mm, inner sep=0.5pt},
  gen/.style={file, fill=ln-accent!15},
  tool/.style={draw, rounded corners=3pt, fill=ln-accent!40, align=center,
               minimum width=22mm, minimum height=6.5mm},
  exe/.style={draw, thick, rounded corners=2pt, fill=ln-box, align=center,
              minimum width=30mm, minimum height=7mm},
  fn/.style={font=\scriptsize\ttfamily},
  what/.style={font=\tiny\sffamily, text=ln-muted},
  lab/.style={font=\scriptsize\sffamily, text=ln-muted}]
  % English descriptions at \scriptsize: \tiny is unreadable on screen, and
  % the widest ("types, constants", 16.4 mm) still fits the 18.5 mm box.
  \newcommand{\lnFile}[2]{{\scriptsize\ttfamily #1}\\[-0.4ex]{\iflangja{\tiny}{\scriptsize}\sffamily\color{ln-muted}#2}}
  % interface and compiler
  \node[file] (x) at (60,52) {\lnFile{sum.x}{\iflangja{XDR インタフェース}{XDR interface}}};
  \node[tool] (rpcgen) at (60,40) {\texttt{rpcgen}};
  \draw[->] (x) -- (rpcgen);
  % source files
  \node[file] (cli) at (10,24) {\lnFile{client.c}{\iflangja{呼び出し元}{caller}}};
  \node[gen]  (cln) at (30,24) {\lnFile{sum\_clnt.c}{\iflangja{クライアントスタブ}{client stub}}};
  \node[gen]  (h)   at (50,24) {\lnFile{sum.h}{\iflangja{型と定数}{types, constants}}};
  \node[gen]  (xdr) at (70,24) {\lnFile{sum\_xdr.c}{\iflangja{XDR ルーチン}{XDR routines}}};
  \node[gen]  (svc) at (90,24) {\lnFile{sum\_svc.c}{\iflangja{サーバスタブ}{server stub}}};
  \node[file] (srv) at (110,24) {\lnFile{server.c}{\iflangja{手続きの実装}{procedure}}};
  \foreach \f in {cln, h, xdr, svc} \draw[->] (rpcgen.south) -- (\f.north);
  % legend
  \node[file, minimum width=5mm, minimum height=3.5mm] (l1) at (3,53) {};
  \node[lab, anchor=west] at (6,53) {\iflangja{開発者が記述}{written by the developer}};
  \node[gen, minimum width=5mm, minimum height=3.5mm] (l2) at (3,48) {};
  \node[lab, anchor=west] at (6,48) {\iflangja{\texttt{rpcgen} が生成}{generated by \texttt{rpcgen}}};
  % executables
  \node[exe] (ce) at (30,5) {\iflangja{クライアントプログラム}{client program}};
  \node[exe] (se) at (90,5) {\iflangja{サーバプログラム}{server program}};
  \foreach \f in {cli, cln, xdr} \draw[->] (\f.south) -- (ce.north);
  \foreach \f in {xdr, svc, srv} \draw[->] (\f.south) -- (se.north);
  \node[lab] at (60,12.5) {\iflangja{コンパイル・リンク}{compile and link}};
\end{tikzpicture}
